\subsection{Interpretation}\label{sec:hazard-interp}

Table~\ref{tab:estimators} summarizes the three estimators evaluated in this paper against the common \$764.7 billion empirical benchmark.

\begin{table}[!tp]
\centering
\footnotesize
\begin{threeparttable}
\caption{Estimator comparison against the \$764.7 billion empirical trapped-liquidity benchmark.}
\label{tab:estimators}
\renewcommand{\arraystretch}{0.92}
\setlength{\itemsep}{0pt}
\begin{tabular}{@{}>{\raggedright\arraybackslash}p{3.2cm}>{\raggedleft\arraybackslash}p{2.4cm}>{\raggedleft\arraybackslash}p{2.4cm}>{\raggedright\arraybackslash}p{2.3cm}>{\raggedright\arraybackslash}p{4.0cm}@{}}
\toprule
Estimator & Trapped liquidity & \% of benchmark & CPR path corr. (lag) & Path diagnostics \\
\midrule
Empirical benchmark (SOMA, phased cap, 30+15yr) & \$764.7B & 100.0\% & -- & mean CPR 5.14\% (ABM-basis back-out; hazard-basis 5.52\%, 5.79\% at the note-rate WAC) \\
Production ABM (Section~\ref{sec:abm-results}) & \$103.7B seed mean (frozen draw \$91.0B) & 13.6\% (seed band 10--17\%; frozen draw 11.9\%) & $r = -0.318$ (0) & mean CPR 11.68\%; raw $R^2$ = $-6.984$ \\
Hazard Path B: literature microsimulation & \$818.5B & 107.0\% (band 105.9--108.2\%) & $r = +0.404$ ($-3$) & mean CPR 4.76\%; $r = +0.190$ at lag 0; 3-month offset (empirical leads; reference only, Section~\ref{sec:pathb}) \\
Hazard Path B: concave-gap variant (Section~\ref{sec:pathb}) & \$808.4B & 105.7\% & $r = +0.281$ (0) & piecewise transform, kink 200 bp; robustness on the log-linear extrapolation \\
Hazard Path A: stratum-month Poisson PML (spec v4, calendar-month) & \$928.9B & 121.5\% & $r = -0.438$ ($-2$) & mean CPR 3.34\%; holdout diagnostics under spec v3 \\
Cross-design ABM, real covariates, recalibrated (Section~\ref{sec:robustness-crossdesign}) & \$453.5B & 59.3\% (fifty-seed mean 60.2\%) & $r = -0.336$ (0) & mean CPR 8.14\%; no positive alignment at any lag \\
Cross-design ABM, real covariates, frozen calibration (Section~\ref{sec:robustness-crossdesign}) & \$159.5B & 20.9\% & $r = -0.311$ (0) & mean CPR 11.63\%; no positive alignment at any lag \\
Hazard Path B: Fannie Mae external replication (Section~\ref{sec:pathb}) & \$821.0B & 107.4\% & -- & null 98.7\%; marginal $+8.68$pp \\
\bottomrule
\end{tabular}
\end{threeparttable}
\end{table}

\noindent{\footnotesize \emph{Notes to Table~\ref{tab:estimators}.} The table reports the three primary estimators, the two cross-design variants of Section~\ref{sec:robustness-crossdesign}, and the Fannie Mae external replication of Path B, whose path-correlation cell is dashed because the committed replication artifact carries headline statistics only (Section~\ref{sec:pathb}). The Danish counterfactual specifications are collected in Table~\ref{tab:danish}. Dashes denote quantities that are not applicable to the row. The benchmark row's ``30+15yr'' scope is not symmetric across estimators: 15-year holdings enter the empirical back-out and the ABM through real weights and real scheduled amortization, with voluntary CPR for those cohorts drawn from the same-coupon 30-year surface (the structural-only fold-in of Appendix~\ref{sec:robustness-pipeline}). Every production ABM figure in this row --- dollars, share, path correlation, $R^2$ and mean CPR --- is that fold-in run, while Section~\ref{sec:robustness-scale}'s scale test sits on the native-gate side of the offset between the two and is reconciled against this row there. The hazard rows are 30-year-collateral estimates from which 15-year face is excluded (Table~\ref{tab:composition}, Term row). Peak-lag entries select the maximum-magnitude cross-correlation (max $|r|$) within a $\pm$3-month window. Under a most-positive-$r$ convention the production ABM's best alignment is $+0.19$ at lag $-3$ (Section~\ref{sec:abm-results}); the two conventions summarize the same lag profile, and Path B's lag $-3$ peak satisfies both. \quad Path B is the only estimator whose aggregate recovery is close to the benchmark-consistent basis (97.9\%, and 100.0\% composed as committed --- 98.1\% converted to a single coupon convention --- at the in-sample calibration point; 91.3\% at the off-window floor the paper headlines). Its raw lag-0 correlation ($+0.190$) is carried by the window's shared trend: in first differences ($-0.23$) and after linear detrending ($-0.20$) the simulated--empirical co-movement is indistinguishable from zero ($\pm 0.31$ at $n = 41$), as it is for every estimator. The two cross-design rows substitute real Freddie Mac structural covariates into the ABM under two calibration choices. Neither reproduces the empirical path's timing --- no estimator does, as above --- but the recalibrated variant's dollar recovery exceeds my own ex-ante threshold for undercutting the paradigm interpretation (Sections~\ref{sec:robustness-crossdesign} and~\ref{sec:conclusion}). Moving-block bootstrap 95\% intervals (block length 6): ABM lag-0 $[-0.60, +0.04]$; Path A lag-0 $[-0.68, -0.02]$; Path B lag-0 $[-0.13, +0.44]$, peak lag $[-0.21, +0.49]$ (no positive lag-0 coefficient is distinguishable from zero; block-length sensitivity at 4/6/8 changes no zero-inclusion verdict in this table except Path A's marginally negative levels endpoint, detailed in Section~\ref{sec:pathb}). Path A's contemporaneous levels correlation is negative on the levels construction, which strengthens the offsetting-errors reading of its aggregate fit (detrended, all three estimators cluster at $-0.20$ to $-0.30$; Section~\ref{sec:pathb}). No timing evidence is claimed for any estimator. Table~\ref{tab:theil} decomposes each estimator's monthly error into bias, variance, and covariance shares and reports cumulative dollar-error profiles; it is a fit diagnostic and changes no verdict in this table. \quad Correlation-column convention: the ABM row's correlations are computed against the ABM pipeline's own empirical back-out (cohort-weighted, term-aware scheduled amortization), while the hazard rows use the hazard-side back-out. Against the common hazard-side series the ABM lag-0 value is $-0.316$, a third-decimal, no-verdict difference. Section~\ref{sec:method-benchmark} notes the two constructions. \quad Shared-accounting-basis equivalents, together with the full-book and composed bases, are collected in Table~\ref{tab:bases} (basis defined in Section~\ref{sec:method-benchmark}, derived in Section~\ref{sec:robustness-hybrid}; the common curtailment netting is a flat 9.1pp for every U.S. hazard leg). Path A's point estimate carries no useful forward precision and is excluded from headline figures. Section~\ref{sec:patha} reports its refit-and-resimulate interval --- since rerun jointly at both specifications, with the joint interval floor-censored and reported beside, not in place of, the committed one --- and the spec v3 provenance of every Path A inferential object.\par}

\begin{table}[H]
\centering
\footnotesize
\begin{threeparttable}
\caption{Recovery of the \$764.748 billion benchmark by accounting basis. Every U.S. hazard leg is quoted on one of these bases: the shared basis nets the common \$69.56 billion curtailment flow from the standalone-scorer figure (a flat 9.1 percentage points), and the composed basis applies the same netting to the book-coupon-reweighted (full-book) figure. The lock-in marginal (central minus $\beta_1 = 0$ null) is basis-invariant because the netting cancels.}
\label{tab:bases}
\begin{tabular}{@{}lrrrrr@{}}
\toprule
Hazard leg & Standalone & Shared & Full-book & Composed & Cum.\ runoff \\
 & (scorer) & ($-9.1$pp) & (book WAC) & ($-9.1$pp) & error (\$B) \\
\midrule
Path B, central (in-sample calibration) & 107.0\% & 97.9\% & 109.1\% & 100.0\% & $+53.8$ \\
Path B, $\beta_1 = 0$ null (in-sample calibration) & 97.8\% & 88.7\% & -- & -- & $-16.6$ \\
\addlinespace
Path B, central (off-window floor) & 100.4\% & 91.3\% & -- & -- & $+2.8$ \\
Path B, $\beta_1 = 0$ null (off-window floor) & 94.8\% & 85.7\% & -- & -- & $-39.8$ \\
\addlinespace
Path A (spec v4) & 121.5\% & 112.4\% & 127.5\% & 118.4\% & $+164.1$ \\
Path B, concave-gap & 105.7\% & 96.6\% & -- & -- & $+43.7$ \\
\bottomrule
\end{tabular}
\end{threeparttable}
\end{table}

\noindent{\footnotesize \emph{Notes to Table~\ref{tab:bases}.} The book-WAC and composed columns are the committed mixed-basis measurements (note rates bucketed against pass-through coupons). Converted to a single pass-through convention they read 107.2\% and 98.1\% for Path B central, with the identified marginal bit-invariant (run \texttt{coupon\_convention\_reweight}). \quad A prior fact frames the whole table. The engine renormalises the simulated pool to realized Fed holdings every month. A leg's dollar roll-off is therefore a \emph{rate} applied to an exogenous level, and its balance path cannot drift from the realized one. No recovery percentage here is a balance-path fit, and none can compound error across months. Each is close to a restatement of that leg's window-mean CPR against the empirical 5.14\% (ABM basis). That is why the sampling intervals are as tight as they are. It is also a limit on what any of these levels can corroborate. The final column is each leg's terminal cumulative runoff error, realized minus simulated roll-off (equivalently, simulated minus realized trapped liquidity) --- the sign-transparent statement of the same information. Recovery is an affine transform of it, $\text{recovery} = 1 + \text{error}/764.7$. So a large additive constant sits inside every percentage in the first four columns, and makes all of them look better. Two things the percentages hide. Path~B's in-sample 107.0\% is an \emph{over}-prediction of trapped liquidity of \$53.8 billion. That is an 8.2\% \emph{under}-prediction of the \$652.8 billion of runoff actually realized---the referent flips even though the map does not. And the headline off-window central leg, at 91.3\% shared, carries the smallest error of any leg here on the standalone scorer: $+\$2.8$ billion, 0.4\% of realized runoff. Errors as a share of realized runoff: $+8.2\%$, $-2.5\%$, $+0.4\%$, $-6.1\%$, $+25.1\%$, $+6.7\%$ down the rows. The lock-in marginal is unaffected by any of this. It is a difference of two trapped aggregates from which the constant cancels --- the same fact as its basis-invariance. Standalone and shared are related by the flat 9.1-percentage-point curtailment wedge of Section~\ref{sec:method-benchmark} (\$69.56 billion over the window), derived in Section~\ref{sec:robustness-hybrid}. Path B's central elasticity band is 105.9--108.2\% standalone (96.8--99.1\% shared); the covariate-prior band is 92.2\%--99.2\% shared (approximately 94.3\%--101.3\% composed). The lock-in marginal is $+9.2$ points ($+\$70.3$ billion) on the U.S. leg at the in-sample calibration point, and $+5.6$ points ($+\$42.6$ billion) at the off-window floor (Section~\ref{sec:robustness-floor}). Both are basis-invariant because the netting is common to the central and null legs. The off-window rows carry the same flat netting. The layer is a constant curtailment flow over a constant cap benchmark, identical across all five committed estimators spanning 29.6 points of standalone recovery, so it is path-invariant and therefore floor-invariant (\texttt{calibration\_reconciliation}). Full-book and composed columns were run only at the in-sample calibration and are not extrapolated to the off-window floor. Path A, the concave-gap variant, the elasticity band, and the covariate-prior band are all in-sample-calibration quantities. The Danish leg's counterfactual image runs at the production floor, so it is comparable to the in-sample point. It is $+\$61.2$ billion, or 8.0 points, the difference reflecting the Danish leg's dynamic-balance compounding (Section~\ref{sec:abm-danish}). The ABM rows of Table~\ref{tab:estimators} run through the accounting layer natively (Figure~\ref{fig:recovery}) and are not restated here.\par}

The two hazard paths share no household-level decision machinery with the ABM. They land closer to the empirical benchmark than any synthetic-population ABM specification tested here (the falsified burnout attempt included) at 97.9\% and 112.4\% in-sample on Section~\ref{sec:robustness-hybrid}'s benchmark-consistent shared accounting basis (107.0\% and 121.5\% under their standalone scorers). At the off-window floor the paper headlines, Path B is 91.3\% shared; Path A was not re-run there. Overshooting less is not validation. Path A overshoots by \$164.1 billion (its terminal cumulative runoff error, a 25.1\% under-prediction of the \$652.8 billion realized) with a monthly path fit no better than the ABM's. Path B's recovery across its full elasticity sensitivity band (\$809.9--\$827.7 billion, 105.9\%--108.2\%; Section~\ref{sec:pathb}) rests on an elasticity imported unmodified from \citet{liebersohn2024}, its verification content limited by Section~\ref{sec:pathb} to the marginal's bounded interval. My claim is therefore narrower than ``the hazard framework fits'' and narrower again than ``lands within 2.1\% of the benchmark'': that 2.1-point miss belongs to the in-sample calibration point, so quoting it beside a marginal measured at the off-window floor mixes two calibrations. The 2.1-point miss is 97.9\%, its composed companion basis-dependent (100.0\% as committed on the mixed coupon basis, 98.1\% on a single pass-through convention, Section~\ref{sec:robustness-hybrid}), retained as the model-fit diagnostic it is, ex-ante covariate priors spanning 92.2\%--99.2\% and approximately 94.3--101.3\% composed, none re-run off-window. At the headline calibration Path B (loan-level survival, no household-utility content) recovers 91.3\% of the aggregate benchmark on the benchmark-consistent accounting basis, a miss of 8.7 points, and 88.2--93.7\% across that floor's defensible range, a miss of 6.3 to 11.8 points. What survives is order-of-magnitude comparison, not precision. Path A lands at 112.4\% sharing data but no estimated parameters, which I read as order-of-magnitude agreement and nothing stronger. It has no useful propagated precision in forward simulation, and a rate-gap coefficient not statistically distinguishable from zero under any bias-respecting construction (Section~\ref{sec:patha}). The same resampling pass puts its mean at 78.3\% of the benchmark with a percentile interval spanning zero, corroborating no sign. No synthetic-population household-choice specification I tested lands within 45\% of it. I attach no timing credential. Path B's raw $+0.190$ lag-0 correlation is carried by the window's shared trend, the co-movement indistinguishable from zero in first differences and after detrending (Section~\ref{sec:pathb}). Table~\ref{tab:estimators}'s moving-block intervals caution against weighing any single timing coefficient, Path B's lag $-3$ peak (interval spanning zero) included. The benchmark side is itself disqualified, its four exact-zero months clip-induced boundary artifacts of the reporting series (Appendix~\ref{sec:robustness-seasonalfloor}, run \texttt{h1\_zero\_months\_diagnosis}). Timing alone does not identify the lock-in channel. The $\beta_1 = 0$ null also peaks at lag $-3$, recovering 88.7\% of the benchmark on the shared basis in-sample and 85.7\% at the off-window floor (Table~\ref{tab:bases}), both under the production floor form. It recovers 35.6\% at that floor under the additive form (Section~\ref{sec:robustness-floor}). The elasticity's identified contribution is its marginal recovery over that null, defined and bounded once in Section~\ref{sec:identification}. The design identifies the marginal, not the lag structure. Across 10 paired in-sample reruns with different draw seeds (central and null sharing each seed) that marginal is \$70.3 billion, varying by less than \$0.01 billion: the engine drains prepayments fractionally and the seed touches only the small delinquency channel. The marginal's uncertainty is therefore the calibration band and structural assumptions, not simulation noise. I restrict this to synthetic-population specifications: Section~\ref{sec:robustness-crossdesign}'s real-covariate cross-design ABM variant does land within 45\% (59.3\%, and 76.3\% once that covariate population is reweighted to the book's own coupon $\times$ vintage composition, within 24\%), complicating the paradigm interpretation this paragraph would otherwise support---and the reweighted reading makes the complication larger, not smaller.

The hazard framework's only behavioral content is loan-level covariates, a cohort-level burnout term, and literature-calibrated elasticities---no analogue to mobility desire, loss aversion, or utility maximization. Its comparative success might therefore be read as evidence for the loan-level and institutional mechanisms Section~\ref{sec:abm}'s synthetic-population ABM does not represent: FICO, LTV, and geography heterogeneity; servicer forbearance and pipeline delay; cohort-level burnout. The paper's own measurements cap that reading. The named mechanisms ablate to a point or less each: burnout 0.86 points at the in-sample floor and 0.45 at the off-window floor, FICO/LTV 1.3 points, the delinquency pipeline \$0.39 billion, with no servicer channel beyond it modeled. Section~\ref{sec:robustness-synthesis}'s synthetic companion, meanwhile, recovers a comparable aggregate level with no real loan-level structure at all. And Section~\ref{sec:robustness-scale}'s data-source asymmetry remains unresolved. The level therefore supports the rate-inelastic decomposition, those mechanisms bounded small and candidates only for the residual. Section~\ref{sec:robustness-scale}'s scale-sensitivity replication rules out the ABM/hazard sample-size difference as the driver of this gap. Section~\ref{sec:robustness-crossdesign}'s cross-design test addresses the data-source difference only partially: real structural covariates there change the ABM's explanatory share substantially. Under one calibration choice that shift is enough to raise the question of whether ``cannot represent'' is the right description at all. The rate-gap coefficient's sign stability across both hazard specifications does not indicate that the underlying par-payoff penalty is doing real work. The reason bears restating: only Path A estimates that coefficient. Path B imports it from \citet{liebersohn2024} as a signed elasticity band that cannot come out the other way. Their agreement, having no freedom to disagree, is an orientation check on unit and sign conventions, not corroboration. It was worth running---it would catch a transposed sign between the paths, and it passes. But the evidence that the penalty has a real effect is the external literature. Path A's estimated coefficient is directionally consistent without being statistically distinguishable from zero under bias-respecting constructions (Section~\ref{sec:patha}). Household rationality is not irrelevant. But the aggregation structure of loan-level survival is the minimum a model of this shortfall must get right, and Section~\ref{sec:robustness-crossdesign} leaves open how much more. That sharpens the two-sided exchange this paper describes, with the qualification Section~\ref{sec:robustness-crossdesign} adds: the price of the TBA market's homogeneity appears to be paid at least in part through loan-level and institutional frictions that the synthetic-population household model does not represent---frictions a market-value repurchase design would substantially relieve, as Section~\ref{sec:abm-danish}'s counterfactual suggests.
\clearpage